
\documentclass[11pt,a4paper]{article}

\usepackage[utf8]{inputenc}
\usepackage[T1]{fontenc}
\usepackage[french]{babel}
\usepackage{lmodern}
\usepackage{microtype}
\usepackage[a4paper,margin=1.65cm,headheight=15pt]{geometry}
\usepackage{amsmath,amssymb,mathtools}
\usepackage{array,tabularx,multicol}
\usepackage{enumitem}
\usepackage[most]{tcolorbox}
\usepackage{xcolor}
\usepackage{tikz}
\usetikzlibrary{arrows.meta,calc}
\usepackage{pdflscape}
\usepackage{fancyhdr}
\usepackage{fancyvrb}
\usepackage{listings}
\usepackage{hyperref}

\definecolor{fondpage}{RGB}{247,248,250}
\definecolor{encre}{RGB}{31,35,40}
\definecolor{bleu}{RGB}{40,91,135}
\definecolor{vert}{RGB}{55,125,92}
\definecolor{orange}{RGB}{178,105,42}
\definecolor{violet}{RGB}{101,77,145}
\definecolor{rouge}{RGB}{170,72,72}
\definecolor{gris}{RGB}{86,91,99}
\definecolor{bleufonce}{RGB}{28,55,120}
\definecolor{bleuclair}{RGB}{238,243,255}
\definecolor{grisclair}{RGB}{245,245,245}
\definecolor{tableline}{RGB}{45,45,45}
\definecolor{softgray}{RGB}{246,247,249}
\definecolor{deepblue}{RGB}{25,62,115}
\definecolor{softblue}{RGB}{235,241,250}
\definecolor{accent}{RGB}{20,82,130}

\hypersetup{colorlinks=true,linkcolor=bleufonce,urlcolor=bleufonce,pdftitle={Parcours de révision -- Suites -- Terminale spécialité},pdfauthor={}}

\setlength{\parindent}{0pt}
\setlength{\parskip}{0.25em}
\renewcommand{\arraystretch}{1.12}
\setlength{\tabcolsep}{3pt}
\setlist[itemize]{leftmargin=1.25em,itemsep=0.15em,topsep=0.2em}
\setlist[enumerate,1]{label=\textbf{\arabic*.}, leftmargin=1.05cm, itemsep=0.3em, topsep=0.2em}
\setlist[enumerate,2]{label=\textbf{\alph*.}, leftmargin=0.85cm, itemsep=0.2em, topsep=0.15em}

\pagestyle{fancy}
\fancyhf{}
\lhead{}
\rhead{}
\cfoot{\thepage}

\newcommand{\R}{\mathbb{R}}
\newcommand{\N}{\mathbb{N}}
\newcommand{\Un}{\ensuremath{(u_n)}}
\newcommand{\e}{\mathrm e}
\newcommand{\vect}[1]{\overrightarrow{#1}}
\newcommand{\origine}[1]{ {\scriptsize\textcolor{bleufonce}{(site #1)}}}

\newcounter{question}
\newcounter{reponse}
\newenvironment{blocquestions}[1]{%
  \begin{tcolorbox}[enhanced,breakable,colback=bleuclair,colframe=bleufonce,boxrule=0.6pt,arc=2mm,left=2mm,right=2mm,top=2mm,bottom=1.5mm,title={\bfseries #1},coltitle=white,colbacktitle=bleufonce,before skip=3mm,after skip=3mm, fontupper=\large]
  \large
  \begin{enumerate}[leftmargin=*,label=\textbf{\arabic*.},itemsep=0.85em,topsep=0.5em]
  \setcounter{enumi}{\value{question}}
}{%
  \setcounter{question}{\value{enumi}}
  \end{enumerate}
  \end{tcolorbox}
}
\newenvironment{blocreponses}[1]{%
  \begin{tcolorbox}[enhanced,breakable,colback=bleuclair,colframe=bleufonce,boxrule=0.6pt,arc=2mm,left=2mm,right=2mm,top=2mm,bottom=1.5mm,title={\bfseries #1},coltitle=white,colbacktitle=bleufonce,before skip=3mm,after skip=3mm, fontupper=\large]
  \large
  \begin{enumerate}[leftmargin=*,label=\textbf{\arabic*.},itemsep=0.45em,topsep=0.5em]
  \setcounter{enumi}{\value{reponse}}
}{%
  \setcounter{reponse}{\value{enumi}}
  \end{enumerate}
  \end{tcolorbox}
}

\newtcolorbox{correctionbox}[1]{enhanced,breakable,colback=grisclair,colframe=black!55,boxrule=0.55pt,arc=2mm,left=2mm,right=2mm,top=2mm,bottom=1.5mm,title=\textbf{#1},coltitle=black,colbacktitle=black!12,before skip=2mm,after skip=2mm}

\newcommand{\titreSujet}[2]{%
  \subsection*{#1}
  \addcontentsline{toc}{subsection}{#1}
  \textit{Référence : #2}\par\medskip\hrule\medskip
}

\newcommand{\sep}{\vspace{0.55mm}}
\tcbset{enhanced,arc=2mm,outer arc=2mm,boxrule=0.42pt,left=1.15mm,right=1.15mm,top=0.75mm,bottom=0.75mm,boxsep=0.35mm,before skip=0.8mm,after skip=0.8mm,fonttitle=\bfseries\footnotesize,coltitle=encre,before upper=\raggedright\fontsize{7.65}{8.15}\selectfont}
\newtcolorbox{fichebox}[3][]{colback=white,colframe=#2!72!black,colbacktitle=#2!18,coltitle=#2!50!black,borderline west={0.9mm}{0pt}{#2!78!black},title={#3},drop fuzzy shadow=black!5,#1}
\newcommand{\tagcol}[2]{\begin{tcolorbox}[enhanced,colback=#1!11,colframe=#1!45!black,boxrule=0.42pt,arc=2mm,left=1mm,right=1mm,top=0.45mm,bottom=0.45mm,before skip=0mm,after skip=0.85mm,fontupper=\bfseries\scriptsize,colupper=#1!55!black]\centering #2\end{tcolorbox}}
\newtcolorbox{panel}[1]{colback=fondpage,colframe=fondpage,boxrule=0pt,arc=0mm,left=0mm,right=0mm,top=0mm,bottom=0mm,height=168mm,valign=top,before skip=0mm,after skip=0mm}
\newcommand{\titrepage}{\begin{tcolorbox}[colback=white,colframe=bleu!70!black,boxrule=0.65pt,arc=3mm,left=2.5mm,right=2.5mm,top=1.15mm,bottom=1.15mm,before skip=0mm,after skip=1.4mm,drop fuzzy shadow=black!10]
{\Large\bfseries Parcours de révision -- Suites}\end{tcolorbox}}
\tikzset{tabouter/.style={draw=tableline, line width=0.85pt},tabline/.style={draw=tableline, line width=0.55pt},forbidden/.style={draw=tableline, line width=0.95pt},vararrow/.style={-{Latex[length=2.7mm,width=2.0mm]}, line width=0.85pt, draw=accent},tabtext/.style={font=\small},tabmath/.style={font=\small},labelcell/.style={font=\small\bfseries, text=deepblue},pointvalue/.style={font=\small\bfseries},endvalue/.style={font=\small}}
\lstset{language=Python,basicstyle=\ttfamily\small,numbers=left,numberstyle=\tiny,frame=single,showstringspaces=false,columns=fullflexible,keepspaces=true}

\begin{document}
\begin{center}
{\LARGE\bfseries Parcours de révision -- Suites}\\[1mm]

\end{center}
\vspace{2mm}
\tableofcontents
\clearpage

\phantomsection
\addcontentsline{toc}{section}{Partie 1 -- Récapitulatif de cours}
\newgeometry{margin=6.5mm}
\pagestyle{empty}
\begin{landscape}
\pagecolor{fondpage}\color{encre}
\thispagestyle{empty}

\fontsize{7.85}{8.45}\selectfont

\titrepage

\noindent
% ============================================================
% COLONNE 1
% ============================================================
\begin{minipage}[t]{0.242\linewidth}
\tagcol{bleu}{1. DÉFINIR, CALCULER, VARIER}
\begin{panel}{bleu}

\begin{fichebox}{bleu}{Suite et notation}
Une suite \Un{} associe à chaque rang $n$ un nombre $u_n$.

\medskip
\textbf{À écrire :} $u_n$ est le terme de rang $n$ ; \Un{} désigne la suite entière.

\medskip
Deux modes de définition fréquents :
\[
 u_n=f(n)
 \qquad\text{ou}\qquad
 u_{n+1}=f(u_n).
\]
\end{fichebox}

\begin{fichebox}{orange}{Premiers termes}
\textbf{Suite explicite :} on remplace $n$ par $0$, $1$, $2$, etc.

\medskip
\textbf{Suite par récurrence :} on part du premier terme puis on calcule successivement :
\[
 u_1=f(u_0),\qquad u_2=f(u_1),\qquad u_3=f(u_2).
\]

\medskip
Chaque remplacement doit apparaître dans la copie.
\end{fichebox}

\begin{fichebox}{vert}{Sens de variation}
Pour étudier les variations de \Un, on peut :
\begin{itemize}
    \item étudier le signe de $u_{n+1}-u_n$ ;
    \item si $u_n>0$, comparer $\dfrac{u_{n+1}}{u_n}$ à $1$ ;
    \item si $u_n=f(n)$, utiliser les variations de $f$.
\end{itemize}

\[
\begin{array}{c|c}
 u_{n+1}-u_n\geq 0 & \Un{} \text{ croissante}\\
 u_{n+1}-u_n\leq 0 & \Un{} \text{ décroissante}
\end{array}
\]
\end{fichebox}

\begin{fichebox}{rouge}{Pièges}
\begin{itemize}
    \item Quelques valeurs ne prouvent pas une variation : elles permettent seulement de conjecturer.
    \item Une suite peut être ni croissante ni décroissante.
    \item Pour utiliser le quotient, vérifier d'abord que les termes sont strictement positifs.
\end{itemize}
\end{fichebox}

\end{panel}
\end{minipage}
\hfill
% ============================================================
% COLONNE 2
% ============================================================
\begin{minipage}[t]{0.242\linewidth}
\tagcol{vert}{2. SUITES ARITHMÉTIQUES ET GÉOMÉTRIQUES}
\begin{panel}{vert}

\begin{fichebox}{vert}{Suite arithmétique}
\Un{} est arithmétique de raison $r$ si :
\[
 u_{n+1}=u_n+r.
\]
Pour $n\geq p$ :
\[
 u_n=u_p+(n-p)r.
\]
\textbf{Montrer :} calculer $u_{n+1}-u_n$ ; cette différence doit être constante.

\textbf{Variation :} le signe de $r$ donne le sens de variation.
\end{fichebox}

\begin{fichebox}{vert}{Suite géométrique}
\Un{} est géométrique de raison $q$ si :
\[
 u_{n+1}=q u_n.
\]
Pour $n\geq p$ :
\[
 u_n=u_p q^{\,n-p}.
\]
\textbf{Montrer :} exprimer $u_{n+1}$ en fonction de $u_n$ et obtenir $u_{n+1}=q u_n$.
\end{fichebox}

\begin{fichebox}{orange}{Sommes utiles}
\[
1+2+\cdots+n=\frac{n(n+1)}2.
\]
\textbf{Nombre de termes de $p$ à $n$ :} $n-p+1$.

\textbf{Arithmétique :}
\[
\sum_{k=p}^{n}u_k=(n-p+1)\frac{u_p+u_n}{2}.
\]

\textbf{Géométrique, si $q\neq1$ :}
\[
\sum_{k=0}^{n}q^k=\frac{1-q^{n+1}}{1-q}
\]
\[
\sum_{k=p}^{n}u_k=u_p\frac{1-q^{n-p+1}}{1-q}.
\]
\end{fichebox}

\end{panel}
\end{minipage}
\hfill
% ============================================================
% COLONNE 3
% ============================================================
\begin{minipage}[t]{0.242\linewidth}
\tagcol{violet}{3. RÉCURRENCE ET CONVERGENCE}
\begin{panel}{violet}

\begin{fichebox}{violet}{Rédiger une récurrence}
Pour démontrer une propriété $P_n$ pour tout $n\geq n_0$ :

\medskip
\textbf{Initialisation :} on vérifie que $P_{n_0}$ est vraie.

\medskip
\textbf{Hérédité :} on suppose que $P_k$ est vraie pour un entier $k\geq n_0$ fixé, et on montre que $P_{k+1}$ est vraie.

\medskip
\textbf{Conclusion :} par récurrence, $P_n$ est vraie pour tout $n\geq n_0$.
\end{fichebox}

\begin{fichebox}{vert}{Suites majorées, minorées, bornées}
\begin{itemize}
    \item \Un{} est majorée si $u_n\leq M$ pour tout $n$.
    \item \Un{} est minorée si $u_n\geq m$ pour tout $n$.
    \item \Un{} est bornée si elle est majorée et minorée.
\end{itemize}
\end{fichebox}

\begin{fichebox}{vert}{Théorèmes de convergence}
\begin{itemize}
    \item croissante et majorée $\Longrightarrow$ convergente ;
    \item décroissante et minorée $\Longrightarrow$ convergente ;
    \item croissante non majorée $\Longrightarrow +\infty$ ;
    \item décroissante non minorée $\Longrightarrow -\infty$.
\end{itemize}
\end{fichebox}

\begin{fichebox}{orange}{Limite et point fixe}
Si $u_{n+1}=f(u_n)$, si $f$ est continue et si $u_n\to \ell$, alors :
\[
\ell=f(\ell).
\]
\end{fichebox}


\end{panel}
\end{minipage}
\hfill
% ============================================================
% COLONNE 4
% ============================================================
\begin{minipage}[t]{0.242\linewidth}
\tagcol{orange}{4. LIMITES ET MÉTHODES DE CALCUL}
\begin{panel}{orange}

\begin{fichebox}{bleu}{Converger, diverger}
\[
\lim_{n\to+\infty}u_n=\ell
\]
signifie que les termes $u_n$ deviennent aussi proches de $\ell$ que l'on veut à partir d'un certain rang.

\medskip
\Un{} est convergente si elle admet une limite finie ; sinon elle est divergente.
\end{fichebox}

\begin{fichebox}{vert}{Suites géométriques $q^n$}
\[
\begin{array}{c|c}
q>1 & q^n\to +\infty\\
q=1 & q^n\to 1\\
-1<q<1 & q^n\to 0\\
q\leq -1 & \text{pas de limite}
\end{array}
\]
\end{fichebox}

\begin{fichebox}{vert}{Comparaison et gendarmes}
\textbf{Comparaison :} si $u_n\geq v_n$ et $v_n\to +\infty$, alors $u_n\to +\infty$.

\medskip
Si $u_n\leq v_n$ et $v_n\to -\infty$, alors $u_n\to -\infty$.

\medskip
\textbf{Gendarmes :} si $u_n\leq v_n\leq w_n$ et si $u_n,w_n\to \ell$, alors $v_n\to \ell$.
\end{fichebox}

\begin{fichebox}{orange}{Formes indéterminées}
À traiter avant de conclure :
\[
+\infty-\infty,
\qquad 0\times\infty,
\qquad \frac{\infty}{\infty},
\qquad \frac{0}{0}.
\]

\medskip
\textbf{Réflexes :} factoriser le terme dominant, développer, mettre au même dénominateur, encadrer.
\end{fichebox}


\end{panel}
\end{minipage}


\end{landscape}
\clearpage
\restoregeometry
\pagecolor{white}\color{black}
\pagestyle{fancy}

\section*{Partie 2 -- Questions flashs}
\addcontentsline{toc}{section}{Partie 2 -- Questions flashs}
\setcounter{question}{0}
\begin{blocquestions}{A. Vocabulaire et formules de base}
  \item Qu'est-ce qu'une suite numérique ? \origine{1}
  \item Comment note-t-on en général le terme de rang \(n\) d'une suite ? \origine{2}
  \item Quelle est la différence entre une définition explicite et une définition par récurrence ? \origine{3}
  \item Quand dit-on qu'une suite est arithmétique ? \origine{4}
  \item Comment s'appelle le réel \(r\) dans une suite arithmétique ? \origine{5}
  \item Quelle est la formule explicite d'une suite arithmétique de premier terme \(u_0\) et de raison \(r\) ? \origine{6}
  \item Quand dit-on qu'une suite est géométrique ? \origine{7}
  \item Comment s'appelle le réel \(q\) dans une suite géométrique ? \origine{8}
  \item Quelle est la formule explicite d'une suite géométrique de premier terme \(u_0\) et de raison \(q\) ? \origine{9}
  \item Comment étudie-t-on en général le sens de variation d'une suite ? \origine{10}
  \item Que montre-t-on pour prouver qu'une suite est croissante ? \origine{11}
  \item Que montre-t-on pour prouver qu'une suite est décroissante ? \origine{12}
  \item Dans quel cas peut-on comparer \(\dfrac{u_{n+1}}{u_n}\) ? \origine{13}
  \item Quelle idée utilise-t-on souvent pour une relation \(u_{n+1}=a u_n+b\) ? \origine{18}
  \item Qu'est-ce qu'un point fixe pour la relation \(u_{n+1}=a u_n+b\) ? \origine{19}
  \item Quel type de boucle Python utilise-t-on souvent pour chercher un seuil ? \origine{20}
\end{blocquestions}

\begin{blocquestions}{B. Limites et convergence}
  \item Que signifie \(\displaystyle \lim_{n\to +\infty} u_n=\ell\) ? \origine{87}
  \item Une suite convergente est-elle toujours bornée ? \origine{88}
  \item Une suite qui tend vers \(+\infty\) est-elle convergente ? \origine{89}
  \item Que vaut \(\displaystyle \lim_{n\to +\infty} 0{,}8^n\) ? \origine{90}
  \item Que vaut \(\displaystyle \lim_{n\to +\infty} 1{,}05^n\) ? \origine{91}
  \item Que vaut \(\displaystyle \lim_{n\to +\infty} \frac{n+1}{n+3}\) ? \origine{92}
  \item Quel est le théorème des gendarmes pour les suites ? \origine{93}
  \item Quand une suite croissante converge-t-elle automatiquement ? \origine{94}
  \item Quand une suite décroissante converge-t-elle automatiquement ? \origine{95}
\end{blocquestions}

\begin{blocquestions}{C. Calculs flash}
  \item Que vaut \(u_5\) si \(u_n=2n+1\) ? \origine{203}
  \item Si \(u_0=7\) et \(u_{n+1}=u_n+2\), que vaut \(u_1\) ? \origine{205}
  \item Quelle suite auxiliaire pense-t-on souvent à poser pour \(u_{n+1}=0{,}8u_n+6\) ? \origine{207}
  \item Quel est le point fixe de la relation \(u_{n+1}=0{,}5u_n+4\) ? \origine{208}
\end{blocquestions}

\begin{blocquestions}{D. Suites bornées, monotones et convergentes}
  \item Qu'appelle-t-on une suite majorée ? \origine{261}
  \item Qu'appelle-t-on une suite minorée ? \origine{262}
  \item Qu'appelle-t-on une suite bornée ? \origine{263}
  \item Qu'appelle-t-on une suite monotone ? \origine{264}
  \item Qu'appelle-t-on une suite constante ? \origine{265}
  \item Qu'appelle-t-on une suite convergente ? \origine{305}
\end{blocquestions}

\begin{blocquestions}{E. Rédaction de la récurrence}
  \item Quels sont les points de rédaction importants pour définir la propriété \(P(n)\) au début d'une récurrence ? \origine{388}
  \item Quels sont les points de rédaction importants pour l'initialisation d'une récurrence ? \origine{389}
  \item Quels sont les points de rédaction importants pour l'hérédité d'une récurrence ? \origine{390}
  \item Quels sont les points de rédaction importants pour la conclusion d'une récurrence ? \origine{391}
\end{blocquestions}

\begin{blocquestions}{F. Un calcul direct pour s'entraîner}
  \item Calculer le terme \(u_{12}\) de la suite arithmétique définie par \(u_0=3\) et de raison \(r=4\). \origine{545}
\end{blocquestions}



\clearpage
\section*{Partie 3 -- Sujets de bac}
\addcontentsline{toc}{section}{Partie 3 -- Sujets de bac}
\titreSujet{Sujet 1 -- Suite auxiliaire et algorithme}{Amérique du Nord, 21 mai 2025, jour 1, exercice 2}

On considère la suite numérique $(u_n)$ définie par son premier terme $u_0=2$ et, pour tout entier naturel $n$, par :
\[
u_{n+1}=\dfrac{2u_n+1}{u_n+2}.
\]
On admet que la suite $(u_n)$ est bien définie.

\begin{enumerate}
  \item Calculer le terme $u_1$.
  \item On définit la suite $(a_n)$ pour tout entier naturel $n$, par :
  \[
a_n=\dfrac{u_n}{u_n-1}.
  \]
  On admet que la suite $(a_n)$ est bien définie.
  \begin{enumerate}
    \item Calculer $a_0$ et $a_1$.
    \item Démontrer que, pour tout entier naturel $n$, $a_{n+1}=3a_n-1$.
    \item Démontrer par récurrence que, pour tout entier naturel $n$ supérieur ou égal à $1$,
    \[
a_n\geqslant 3n-1.
    \]
    \item En déduire la limite de la suite $(a_n)$.
  \end{enumerate}
  \item On souhaite étudier la limite de la suite $(u_n)$.
  \begin{enumerate}
    \item Démontrer que, pour tout entier naturel $n$,
    \[
u_n=\dfrac{a_n}{a_n-1}.
    \]
    \item En déduire la limite de la suite $(u_n)$.
  \end{enumerate}
  \item On admet que la suite $(u_n)$ est décroissante. On considère le programme suivant écrit en langage Python :
\begin{Verbatim}[frame=single, fontsize=\small]
def algo(p):
    u = 2
    n = 0
    while u - 1 > p:
        u = (2*u + 1)/(u + 2)
        n = n + 1
    return (n, u)
\end{Verbatim}
  \begin{enumerate}
    \item Interpréter les valeurs $n$ et $u$ renvoyées par l'appel de la fonction \texttt{algo(p)} dans le contexte de l'exercice.
    \item Donner, sans justifier, la valeur de $n$ pour $p=0{,}001$.
  \end{enumerate}
\end{enumerate}

\newpage
\titreSujet{Sujet 2 -- Deux modèles d'évolution d'une population}{Centres étrangers, 13 juin 2025, sujet 2, exercice 1}

On se propose de comparer l'évolution d'une population animale dans deux milieux distincts $A$ et $B$.
Au 1er janvier 2025, on introduit $6\,000$ individus dans chacun des milieux $A$ et $B$.

\textbf{Partie A}

Dans cette partie, on étudie l'évolution de la population dans le milieu $A$.

On suppose que dans ce milieu, l'évolution de la population est modélisée par une suite géométrique $(u_n)$ de premier terme $u_0=6$ et de raison $0{,}93$.
Pour tout entier naturel $n$, $u_n$ représente la population au 1er janvier de l'année $2025+n$, exprimée en millier d'individus.

\begin{enumerate}
  \item Donner, selon ce modèle, la population au 1er janvier 2026.
  \item Pour tout entier naturel $n$, exprimer $u_n$ en fonction de $n$.
  \item Déterminer la limite de la suite $(u_n)$. Interpréter ce résultat dans le contexte de l'exercice.
\end{enumerate}

\textbf{Partie B}

Dans cette partie, on étudie l'évolution de la population dans le milieu $B$.

On suppose que dans ce milieu, l'évolution de la population est modélisée par la suite $(v_n)$ définie par
\[
v_0=6\qquad\text{et, pour tout entier naturel }n,\qquad v_{n+1}=-0{,}05v_n^2+1{,}1v_n.
\]
Pour tout entier naturel $n$, $v_n$ représente la population au 1er janvier de l'année $2025+n$, exprimée en millier d'individus.

\begin{enumerate}
  \item Donner, selon ce modèle, la population au 1er janvier 2026.
\end{enumerate}
Soit $f$ la fonction définie sur l'intervalle $[0;+\infty[$ par
\[
f(x)=-0{,}05x^2+1{,}1x.
\]
\begin{enumerate}[resume]
  \item Démontrer que la fonction $f$ est croissante sur l'intervalle $[0;11]$.
  \item Démontrer par récurrence que, pour tout entier naturel $n$, on a
  \[
  2\leqslant v_{n+1}\leqslant v_n\leqslant 6.
  \]
  \item En déduire que la suite $(v_n)$ est convergente vers une limite $\ell$.
  \item
  \begin{enumerate}
    \item Justifier que la limite $\ell$ vérifie $f(\ell)=\ell$, puis en déduire la valeur de $\ell$.
    \item Interpréter ce résultat dans le contexte de l'exercice.
  \end{enumerate}
\end{enumerate}

\textbf{Partie C}

Cette partie a pour but de comparer l'évolution de la population dans les deux milieux.

\begin{enumerate}
  \item En résolvant une inéquation, déterminer l'année à partir de laquelle la population du milieu $A$ sera strictement inférieure à $3\,000$ individus.
  \item À l'aide de la calculatrice, déterminer l'année à partir de laquelle la population du milieu $B$ sera strictement inférieure à $3\,000$ individus.
  \item Justifier qu'à partir d'une certaine année, la population du milieu $B$ dépassera la population du milieu $A$.
  \item On considère le programme Python ci-dessous.
  \begin{enumerate}
    \item Recopier et compléter ce programme afin qu'après exécution, il affiche l'année à partir de laquelle la population du milieu $B$ est strictement supérieure à la population du milieu $A$.
    \item Déterminer l'année affichée après exécution du programme.
  \end{enumerate}
\end{enumerate}
\begin{Verbatim}[frame=single, fontsize=\small]
n = 0
u = 6
v = 6
while ... :
    u = ...
    v = ...
    n = n + 1
print(2025 + n)
\end{Verbatim}

\newpage
\titreSujet{Sujet 3 -- Convergence de deux suites vers une même limite}{Métropole, 9 septembre 2025, sujet 1, exercice 3}

Le but de cet exercice est d'étudier les convergences de deux suites vers une même limite.

\textbf{Partie A}

On considère la fonction $f$ définie sur $[2;+\infty[$ par
\[
f(x)=\sqrt{3x-2}.
\]

\begin{enumerate}
  \item Justifier les éléments du tableau de variations ci-dessous :
  \begin{center}
  \begin{tikzpicture}[x=1cm,y=1cm]
    \def\L{2.25}
    \def\W{8.60}
    \def\H{2.70}
    \def\Yx{1.90}

    % Fonds
    \fill[softgray] (0,\Yx) rectangle (\W,\H);
    \fill[softblue] (0,0) rectangle (\L,\Yx);

    % Cadre et séparations
    \draw[tabouter] (0,0) rectangle (\W,\H);
    \draw[tabline] (0,\Yx) -- (\W,\Yx);
    \draw[tabline] (\L,0) -- (\L,\H);

    % Libellés
    \node[labelcell] at ({\L/2},{(\Yx+\H)/2}) {$x$};
    \node[labelcell] at ({\L/2},{\Yx/2}) {$f(x)$};

    % Ligne x
    \node[pointvalue] at ({\L+0.35},{(\Yx+\H)/2}) {$2$};
    \node[tabmath] at ({\W-0.45},{(\Yx+\H)/2}) {$+\infty$};

    % Valeurs et flèche
    \node[pointvalue] at ({\L+0.45},0.35) {$2$};
    \node[endvalue] at ({\W-0.65},1.45) {$+\infty$};
    \draw[vararrow] ({\L+0.90},0.50) -- ({\W-1.10},1.35);
  \end{tikzpicture}
  \end{center}
\end{enumerate}

On admet que la suite $(u_n)$ vérifiant $u_0=6$ et, pour tout entier naturel $n$,
\[
u_{n+1}=f(u_n)
\]
est bien définie.

\begin{enumerate}[resume]
  \item
  \begin{enumerate}
    \item Démontrer par récurrence que, pour tout entier naturel $n$,
    \[
    2\leqslant u_{n+1}\leqslant u_n\leqslant 6.
    \]
    \item En déduire que la suite $(u_n)$ converge.
  \end{enumerate}
  \item On appelle $\ell$ la limite de $(u_n)$. On admet qu'elle est solution de l'équation $f(x)=x$. Déterminer la valeur de $\ell$.
  \item On considère la fonction \texttt{rang} écrite ci-dessous en langage Python. On rappelle que \texttt{sqrt(x)} renvoie la racine carrée du nombre $x$.
\begin{Verbatim}[frame=single, fontsize=\small]
from math import *

def rang(a):
    u = 6
    n = 0
    while u >= a:
        u = sqrt(3*u - 2)
        n = n + 1
    return n
\end{Verbatim}
  \begin{enumerate}
    \item Pourquoi peut-on affirmer que \texttt{rang(2.000001)} renvoie une valeur ?
    \item Pour quelles valeurs du paramètre $a$ l'instruction \texttt{rang(a)} renvoie-t-elle un résultat ?
  \end{enumerate}
\end{enumerate}

\textbf{Partie B}

On admet que la suite $(v_n)$ vérifiant $v_0=6$ et, pour tout entier naturel $n$,
\[
v_{n+1}=3-\dfrac{2}{v_n}
\]
est bien définie.

\begin{enumerate}
  \item Calculer $v_1$.
  \item Pour tout entier naturel $n$, on admet que $v_n\neq 2$ et on pose :
  \[
w_n=\dfrac{v_n-1}{v_n-2}.
  \]
  \begin{enumerate}
    \item Démontrer que la suite $(w_n)$ est géométrique de raison $2$ et préciser son premier terme $w_0$.
    \item On admet que, pour tout entier naturel $n$,
    \[
w_n-1=\dfrac{1}{v_n-2}.
    \]
    En déduire que, pour tout entier naturel $n$,
    \[
v_n=2+\dfrac{1}{1{,}25\times 2^n-1}.
    \]
    \item Calculer la limite de $(v_n)$.
  \end{enumerate}
  \item Déterminer le plus petit entier naturel $n$ pour lequel $v_n<2{,}01$ en résolvant l'inéquation.
\end{enumerate}

\textbf{Partie C}

À l'aide des parties précédentes, déterminer le plus petit entier $N$ tel que, pour tout $n\geqslant N$, les termes $v_n$ et $u_n$ appartiennent à l'intervalle $]1{,}99;2{,}01[$.


\end{document}
